% Section 5: Composite Systems and Local Tomography
% \input-able LaTeX file for Paper 5
% Conventions: a \sp b = sequential product (dot notation matching vdW)
%   C_p = Alfsen-Shultz compression
%   tau = normalized EJA trace
%   phi = tracking isomorphism (self-model)

\section{Composite Systems and Local Tomography}\label{sec:composite}

%--------------------------------------------------------------------
\subsection{The Body--Model Composite}\label{subsec:composite-def}

Having established that the state space $V$ of the body $B$ (and
isomorphically the model $M$) carries EJA structure, we now formalize the
composite system $V_{BM}$ using only generalized probabilistic theory
(GPT) primitives~\cite{Plavala2023,Barnum2023}.

\begin{definition}[Composite OUS]\label{def:composite}
Let $V_B$ and $V_M$ be finite-dimensional order unit spaces, each
equipped with a sequential product satisfying S1--S7. The
\emph{composite order unit space} $V_{BM}$ is defined by:
\begin{enumerate}
\item[\emph{(C1)}] \textbf{Vector space.}
  $V_{BM}$ is a finite-dimensional real vector space containing
  $V_B \otimes V_M$ (the algebraic tensor product of real vector
  spaces) as a subspace.
\item[\emph{(C2)}] \textbf{Order unit.}
  $\mathbf{1}_{BM} = \mathbf{1}_B \otimes \mathbf{1}_M$.
\item[\emph{(C3)}] \textbf{Product states.}
  For every state $\omega_B$ on $V_B$ and $\omega_M$ on $V_M$, the
  product state $\omega_B \otimes \omega_M$ belongs to $\mathcal{S}(V_{BM})$,
  defined by $(\omega_B\otimes\omega_M)(a\otimes b)
  = \omega_B(a)\,\omega_M(b)$.
\item[\emph{(C4)}] \textbf{Non-signaling.}
  For every joint state $\omega \in \mathcal{S}(V_{BM})$, the marginals
  $\omega_B(a) \coloneqq \omega(a\otimes\mathbf{1}_M)$ and
  $\omega_M(b) \coloneqq \omega(\mathbf{1}_B\otimes b)$ are well-defined
  states, independent of what is measured on the other subsystem.
\end{enumerate}
\end{definition}

\begin{assumption}[Minimal composite]\label{ass:minimal-composite}
$V_{BM}$ is the \emph{minimal} order unit space satisfying
(C1)--(C4) and the product-form sequential product
(Definition~\ref{def:product-sp} below).
\end{assumption}

This construction uses only real vector spaces, compact convex state
spaces, affine functionals, and the non-signaling constraint --- no
Hilbert space tensor products, complex numbers, or density matrices.

%--------------------------------------------------------------------
\subsection{Product-Form Sequential Product}\label{subsec:product-sp}

\begin{definition}[Product-form SP]\label{def:product-sp}
The sequential product on $V_{BM}$ is defined on product effects by
\begin{equation}\label{eq:product-sp}
  \sp{(a \otimes b)}{(c \otimes d)}
  = (\sp{a}{c}) \otimes (\sp{b}{d}),
\end{equation}
where $\sp{\cdot}{\cdot}$ on the right-hand side denotes the
factor-level sequential product~\eqref{eq:sp-explicit}.
\end{definition}

\begin{proposition}[S1--S7 inheritance]\label{prop:inheritance}
If the factor sequential products on $V_B$ and $V_M$ each satisfy
S1--S7, then the product-form sequential product on $V_{BM}$ satisfies
S1--S7.
\end{proposition}

\begin{proof}[Proof sketch]
Axioms S1 (additivity), S2 (continuity), and S3 (unitality) follow
directly from the corresponding factor-level axioms and the bilinearity
of the tensor product. Axioms S5--S7 (compatible associativity,
additivity, multiplicativity) reduce to the factor-level axioms because
compatibility of product effects
$a\otimes b$ and $c\otimes d$ follows from factor-wise compatibility
$\sp{a}{c}=\sp{c}{a}$ and $\sp{b}{d}=\sp{d}{b}$.

The inheritance of S4 (orthogonality symmetry) is more subtle. If
$\sp{(a\otimes b)}{(c\otimes d)} = 0$, then
$(\sp{a}{c})\otimes(\sp{b}{d}) = 0$. In an order unit space, states
separate points (Alfsen--Shultz~\cite{AlfsenShultz2003}, Theorem~1.23):
for product states $\omega_B\otimes\omega_M$, the evaluation
$\omega_B(\sp{a}{c})\,\omega_M(\sp{b}{d}) = 0$ for all
$\omega_B,\omega_M$ forces either $\sp{a}{c} = 0$ or $\sp{b}{d} = 0$.
By the factor-level S4, the corresponding reverse products vanish, giving
$\sp{(c\otimes d)}{(a\otimes b)} = 0$.
\end{proof}

\begin{remark}[Bootstrapping order]\label{rem:bootstrap}
The verification of S1--S7 on the composite avoids a potential
circularity.  The product-form sequential
product~\eqref{eq:product-sp} is first defined on the algebraic
tensor product $V_B \otimes V_M$ (not on all of~$V_{BM}$) by
bilinear extension.  Axioms S1--S7 are verified on this subspace
using only the factor-level axioms and bilinearity---no reference
to the full composite $V_{BM}$ is needed.  By
Theorem~\ref{thm:EJA}, the algebraic tensor product equipped with
this product carries EJA structure.  The trace form non-degeneracy
argument (Proposition~\ref{prop:nondeg}) then shows that the
product effects $\{a_i \otimes b_j\}$ are linearly independent,
giving $\dim(V_{BM}) \ge \dim(V_B) \cdot \dim(V_M)$.  Minimality
(Assumption~\ref{ass:minimal-composite}) supplies the upper bound,
yielding $V_{BM} = V_B \otimes V_M$.
\end{remark}

%--------------------------------------------------------------------
\subsection{From Faithful Tracking to Local
Tomography}\label{subsec:lt-proof}

The key step is showing that faithful self-modeling forces the composite
to be locally tomographic: product measurements suffice to determine all
joint states.

\begin{assumption}[Simple EJA]\label{ass:simple-eja}
The EJA $V$ is \emph{simple} (not a nontrivial direct sum of EJAs).
\end{assumption}

This assumption is needed for the non-degeneracy of the trace form below.

\begin{definition}[Correlation bilinear form]\label{def:corr-form}
Let $\tau$ denote the normalized trace functional on the EJA~$V$
(Faraut--Kor\'anyi~\cite{FarautKoranyi1994}, Ch.~III) and let $\circ$
denote the Jordan product. Define the \emph{correlation bilinear form}
\begin{equation}\label{eq:corr-form}
  B(a,b) \;=\; \tau\!\left(a \circ \varphi^{-1}(b)\right),
\end{equation}
where $\varphi^{-1} \colon V_M \to V_B$ is the inverse of the tracking
isomorphism.
\end{definition}

\begin{proposition}[Non-degeneracy]\label{prop:nondeg}
On a simple EJA, the bilinear form $B(a,b)$ is non-degenerate.
\end{proposition}

\begin{proof}
On a simple EJA, the trace form $(a,c) \mapsto \tau(a\circ c)$ is
non-degenerate (Faraut--Kor\'anyi~\cite{FarautKoranyi1994},
Proposition~III.4.2). Since $\varphi\colon V_B\to V_M$ is an
isomorphism (faithful self-modeling), $\varphi^{-1}$ is a bijection,
so $B(a,b) = \tau(a\circ\varphi^{-1}(b))$ is non-degenerate: if
$B(a,b) = 0$ for all $a$, then $\varphi^{-1}(b) = 0$, hence $b=0$.
\end{proof}

\begin{theorem}[Local tomography]\label{thm:local-tomo}
Under Assumptions~\ref{ass:minimal-composite} and~\ref{ass:simple-eja},
\[
  \dim(V_{BM}) \;=\; \dim(V_B)\cdot\dim(V_M).
\]
\end{theorem}

\begin{proof}[Proof sketch]
The non-degeneracy of $B$ implies that the product effects
$\{a_i\otimes b_j\}$, where $\{a_i\}$ and $\{b_j\}$ are bases for $V_B$
and $V_M$ respectively, are linearly independent in $V_{BM}$. Indeed, if
$\sum_{i,j}\alpha_{ij}\,a_i\otimes b_j = 0$, then for any fixed~$a$,
the map $b \mapsto \sum_j \alpha_{ij}\,B(a,b_j) = 0$ forces all
$\alpha_{ij}=0$ by non-degeneracy.

This gives the lower bound $\dim(V_{BM}) \ge \dim(V_B)\cdot\dim(V_M)$.
By minimality (Assumption~\ref{ass:minimal-composite}), the composite
contains no structure beyond what is forced by axioms (C1)--(C4) and
the product-form SP. Since the product effects span a
$\dim(V_B)\cdot\dim(V_M)$-dimensional subspace, and the SP preserves
this subspace (product effects map to product effects
under~\eqref{eq:product-sp}), minimality gives the upper bound
$\dim(V_{BM}) \le \dim(V_B)\cdot\dim(V_M)$.
\end{proof}

%--------------------------------------------------------------------
\subsection{The Entangled Sector}\label{subsec:entangled}

The \emph{entangled sector} consists of states in $V_{BM}$ not reachable
as convex combinations of product states. Minimality
(Assumption~\ref{ass:minimal-composite}) eliminates this sector by
construction: $V_{BM}$ is defined as the smallest OUS satisfying the
composite axioms, so it contains no ``hidden'' entangled structure beyond
what the sequential product and non-signaling constraints require.

For complex quantum mechanics ($V = M_n(\mathbb{C})^{\mathrm{sa}}$),
the minimal and maximal composites
coincide~\cite{Barnum2023}: the standard Hilbert space
tensor product $M_{nm}(\mathbb{C})^{\mathrm{sa}}$ is the unique composite
satisfying (C1)--(C4). This is why minimality is not a restrictive
assumption for the complex case.

For real or quaternionic quantum mechanics, the minimal and maximal
composites differ. This discrepancy manifests as a failure of local
tomography and is the mechanism by which these types are excluded:

\begin{remark}[Negative checks]\label{rem:negative-checks}
Dimension counting reveals that local tomography fails for all
non-complex types at $n=2$:
\begin{center}
\begin{tabular}{lcccl}
\toprule
Type & $d = \dim(V)$ & $d^2$ & $\dim(\text{composite})$
& Local tomography? \\
\midrule
$M_2(\mathbb{R})^{\mathrm{sa}}$ & $3$ & $9$ & $10$ & No\; ($9 < 10$) \\
$M_2(\mathbb{C})^{\mathrm{sa}}$ & $4$ & $16$ & $16$ & Yes\; ($16 = 16$)\\
$M_2(\mathbb{H})^{\mathrm{sa}}$ & $6$ & $36$ & $28$ & No\; ($36 > 28$)\\
\bottomrule
\end{tabular}
\end{center}
The real composite is ``too large'' (product measurements cannot
distinguish all joint states), while the quaternionic composite is ``too
small'' (product measurements are linearly dependent on the joint system).
Only the complex type achieves the exact match $d^2 =
\dim(\text{composite})$.
\end{remark}
